\documentclass[letterpaper,times]{IONconf-v2}
\addbibresource{srnx-paper.bib}
\usepackage[hidelinks]{hyperref}
\usepackage{bytefield}

%%%%%%%%%%%%%%%%%%%%%%%%%%%%%%%%%%%%%%%%%%%%%%%%%%%%%%%%%%%%%%%%%%%%%%%%%
%% Document starts here

\articletype{Short paper}

%\received{26 April 2016}
%\revised{6 June 2016}
%\accepted{6 June 2016}
%\doi{10.1109/xxx.xx.xxxx}
%\journalname{NAVIGATION}
%\journalvolume{5}
%\journalnumber{7}

% Authors are grouped by institution. Institution names are italicized.
\title{A Space- and Time-Efficient Format for RINEX Data}
\author[1]{Michael Poole}
\authormark{Michael Poole}
\address[1]{ION Member, Virginia, United States}
\corres{Michael Poole\\ \email{mdpoole@troilus.org}}

\begin{document}

\abstract[Abstract]{The Receiver-Independent Exchange (RINEX) data format
is a simple text-oriented format that represents data in an epoch-major
or "row-oriented" format.
It provides a simple and robust format, but requires significant space
to store.
Y.~Hatanaka has defined several iterations of the Compressed RINEX
(CRINEX) format to represent observation data in a compressed text format.
Online archives often use text-to-binary compression formats to further
compress both RINEX and CRINEX files.
This paper describes an observation-major (columnar) binary format that
requires less storage space than these formats while supporting very
fast reads and competitive write times.}

\keywords{RINEX, CRINEX, compression, columnar data, delta coding}

\maketitle

\section{Introduction}

This paper describes the "Succinct RINEX" (SRNX) file format, a
representation of RINEX observation data that leverages techniques that
have been adopted in the "big data" field to improve data density and
access times.

Section 2 briefly describes the file format.
Section 3 provides performance for storage density, conversion times
from RINEX and CRINEX (\cite{RNXCMP}) formats, and time to load data
from files.
Section 4 describes additional features of the reference implementation
of the format and suggests potential further work.

\section{The SRNX Format}

An SRNX file consists of a number of chunks with an optional file-level
digest following the last chunk.
Each chunk starts with a four-character code (FOURCC) identifying its
semantic and structural type, followed by the length of the payload in
little-endian base-128 (LEB128) format, the payload, and the per-chunk
digest.
The per-chunk digest covers the FOURCC, length and payload fields.
Figure \ref{fig:srnx-file} shows the general structure of a SRNX file;
Figure \ref{fig:srnx-chunk} illustrates the FOURCC structure for a
file-initial SRNX chunk.

\begin{figure}
\centering
\begin{minipage}{0.45\textwidth}
	\centering
	\includegraphics[width=\textwidth]{figures/srnx-file.pdf}
	\caption{SRNX File Structure}
	\label{fig:srnx-file}
\end{minipage}\hfill
\begin{minipage}{0.45\textwidth}
	\centering
	\includegraphics[width=\textwidth]{figures/pareto.pdf}
	\caption{Compression Ratio vs Wall Time}
	\label{fig:compress-time}
\end{minipage}
\end{figure}

\begin{figure}[htb]
\centering
\begin{bytefield}[bitwidth=1.1em]{32}
	\bitheader{0-31} \\
	\begin{rightwordgroup}{FOURCC \\ header}
		\bitbox{32}{'SRNX'} \\
		\bitbox{8}{payload length=16}
	\end{rightwordgroup} \\
	\begin{rightwordgroup}{Payload}
		\bitboxes{8}{{major=1}{minor=0}{chunk=2 (CRC-32C)}{file=5 (BLAKE3-256)}} \\
		\bitbox{32}{SDIR offset=27710945 (\texttt{e1 ab 9b 0d})} \\
		\wordbox{2}{padding (zero bytes)}
	\end{rightwordgroup} \\
	\begin{rightwordgroup}{Digest}
		\wordbox{1}{CRC-32C (\texttt{1e 72 4f ba})}
	\end{rightwordgroup}
\end{bytefield}
\caption{SRNX Chunk Example}
\label{fig:srnx-chunk}
\end{figure}

\begin{figure}[htb]
\centering
\begin{bytefield}[bitwidth=1.1em]{32}
	\bitheader{0-31} \\
	\begin{rightwordgroup}{FOURCC \\ header}
		\bitbox{32}{'SOCD'} \\
		\bitbox{24}{payload length}
	\end{rightwordgroup} \\
	\begin{rightwordgroup}{Payload}
		\bitbox{24}{RINEX satellite code (e.g. 'G01')}\bitbox{8}{pad=0} \\
		\bitbox{24}{RINEX observation code (zero padded for v2)}\bitbox{8}{pad=0} \\
		\bitbox{16}{Length of LLI data} \bitbox[lt]{16}{RLE LLIs} \\
		\bitbox{32}{ ... run-length encoded LLIs ... } \\
		\bitbox{16}{Length of SSI data} \bitbox[lt]{16}{RLE SSIs} \\
		\bitbox{32}{ ... run-length encoded SSIs ... } \\
		\bitbox{24}{Length of obs data} \bitbox{8}{Flag+Order} \\
		\bitbox{32}{Optional scale (LEB128)} \\
		\bitbox{32}{\textit{Order} signed LEB128 initial deltas} \\
		\wordbox{2}{... compressed block data ...}
	\end{rightwordgroup} \\
	\begin{rightwordgroup}{Digest}
		\wordbox{1}{CRC-32C}
	\end{rightwordgroup}
\end{bytefield}
\caption{SOCD Chunk Format}
\label{fig:socd-chunk}
\end{figure}

\subsection{Chunk Types}

The initial SRNX chunk identifies the file as a SRNX; it provides a
version, message-digest selection for each chunk and for the file as a
whole, and an offset of the SDIR chunk (if one exists).
CRC-32C is the recommended digest for each chunk; BLAKE3-256 is the
recommended whole-file digest.

The following RHDR chunk carries the RINEX header, with newlines
normalized to Unix style and trailing whitespace removed from lines to
facilitate finding lines by type.

The EPOC chunk describes the timestamps of observation epochs.

Each EVTF chunk contains one special event record, with epoch flag 2
through 5 inclusive, with newlines normalized to Unix style.

Each observed satellite has one or more SOCD chunks that contain packed
observation code data, and typically has an SATE chunk that provides a
directory of SOCD chunks for name-based lookup.
The basic structure of an SOCD chunk is zero-filled satellite and
observation code identifiers, run-length coded LLIs and SSIs, the delta
coding order plus a flag indicating the presence of a scale value,
the optional scale indicating the resolution of the values
(for example, when $C/N_0$ is rounded to quarter-dB-Hz resolution),
initial values for the delta coder, and a number of compressed blocks.
The structure of the SOCD chunk is shown in Figure \ref{fig:socd-chunk}.

Each compressed SOCD block has a header byte indicating the format.
224 values select one of seven matrix lengths, representing the number
of coded values in the matrix, and the bit length of each value from 1
to 32 bits.
Three values select special cases, starting with a LEB128 count of
(respectively) absent values, zero values or signed LEB128 values.
The other 29 header values are reserved.
The compressed block form was inspired by \cite{https://doi.org/10.1002/spe.2203}.

Finally, there is typically an SDIR chunk providing a directory of SATE
chunks to find per-satellite information.

A detailed specification of the format for each chunk type is provided
in the reference implementation.

\subsection{Compression Algorithm(s)}

SRNX uses run-length encoding for several fields: epoch timestamps,
LLIs and SSIs.
It uses a CRINEX-like delta coding for observation values, and codes the
residuals with a variable number of bits over time.

The reference compressor identifies the optimal coded length for a
candidate delta order, and will attempt the next largest order until
either order 7 is checked or two consecutive orders increase the size.
It identifies the optimal coded length for a given order with dynamic
programming to find the shortest path from the start to end of the
observations, with different values of header bytes in SOCD blocks
(plus number of LEB128 values) representing edges in the graph.
This results in a runtime less than twice that of \texttt{RNX2CRX}.

For an arbitrary day in May 2026, the most common delta order was 3,
followed by 1 and 2, with very few SOCDs using order 0 or 4, and none
using more than 4.

\section{Performance}

\subsection{Compression Ratio and Time}

Figure \ref{fig:compress-time} shows the total time to compress a one-day
set of NOAA CORS station data and the corresponding compression ratio
achieved by various schemes, with numbers shown in
Table \ref{tab:compress-time}.
Wall time is measured with parallelism of 16 jobs.
This data is from 9 August 2022 because GLONASS satellite R25 began
transmitting the next day, and the final version of teqc rejects that.

Notably, general-purpose compressors provide only a few percent gain on
an SRNX file, although 30-second decimated files are small enough that
the wall time for further compression may be tolerable for that gain.

\begin{table}[htb]
\caption{Compression Ratio vs Wall Time by compression method}
\label{tab:compress-time}
\begin{tblr}{colspec={X[l]X[c]X[c]X[l]},width=\textwidth,rowsep=2pt,
	row{even} = {white,font=\small},
	row{odd} = {bg=black!10,font=\small},
	row{1} = {bg=black!20,font=\bfseries\small},
	hline{Z} = {1pt,solide,black!60}}
Method & Ratio & Wall Time (s) & Comments \\
CRINEX & 4.1636 & 102.36 & RNXCMP v4.2.0 \\
CRINEX + gzip & 11.7127 & 1039.59 & gzip v1.14, -9 \\
CRINEX + bzip3 & 17.3284 & 660.2 & bzip3 v1.5.3 \\
SRNX & 19.3294 & 175.3 & this work \\
SRNX + bzip3 & 19.6633 & 500.35 & \\
SRNX + zpaq & 19.8464 & 3427.45 & zpaq v7.15, -m56 \\
30 s & 15.8571 & 57.15 & teqc 2019Feb25 \\
30 s + CRINEX & 52.7277 & 64.14 & \\
30 s + CRINEX + gzip & 138.2165 & 125.72 & \\
30 s + SRNX & 214.2017 & 69.91 & \\
30 s + SRNX + pcompress & 223.7749 & 100.18 & pcompress v3.1, -a -l 14 -D -G -L -s 64m \\
30 s + SRNX + zpaq & 226.9678 & 360.94 & \\
\end{tblr}
\end{table}

\subsection{Loading Time}

Reading SRNX data is time-competitive with reading RINEX and CRINEX data.
When comparing mismatched pairs (thus triggering termination soon after
loading the files) on a Macbook~Pro with M2~Max processor, the reference
implementation's \texttt{rnxcmp} utility loads files from the operating
system's cache three to four times as fast as reading content-equivalent
RINEX and CRINEX into the same columnar in-memory format.

(The reference implementation offers a row-based API to read RINEX and
CRINEX, avoiding memory management to put data in column format.
This mode takes 2.5 times as long as reading SRNX.)

\subsection{Correctness}

Underlying data and scripts, plus links to two implementations (a
reference implementation in C and an AI-assisted clean-room
implementation in Rust) are available in the supplementary data
(\cite{REPO}).

Both implementations include semantic comparison scripts, which gave
clean results for 9 August 2022, an earlier benchmark day (18 July 2020),
and most of May 2026.
In May 2026, an average of less than one file per day was rejected by
the \texttt{rnx2srnx} converter due to syntax errors in the original
RINEX files; there were no mismatches among converted files.

\section{Further Work}

Many of the RINEX files rejected by the reference implementation could
be recovered given the redundancies in RINEX; common error modes include
including a new satellite only in the epoch header or only in the
observations, and specifying the satellite code with spaces in place of
the identifying digits.

The seven matrix lengths for SOCD blocks (8, 16, 24, 32, 40, 48 and 112)
were chosen by optimizing compression size over the benchmark day, and
might not be globally optimal.
Varying these per file would be straightforward, although selecting them
is expensive during compression.

Separating fields from an SRNX and trying to compress them individually
identified likely redundancy in the SOCD block-header bytes, and
redundancy in the run-length coded form of SSIs.
LLI and SSI might benefit from using the observation coding form, and
SSI could be elided when it is a function of $C/N_0$.

\printbibliography[title=References]

\end{document}
